\documentclass[12pt,a4paper]{article}
\usepackage[utf8]{inputenc}
\usepackage[T1]{fontenc}
\usepackage[spanish]{babel}
\usepackage{amsmath,amssymb}
\usepackage{geometry}
\usepackage{siunitx}
\usepackage{booktabs}
\usepackage{graphicx}
\geometry{margin=2.5cm}

\title{Tarea --- Cálculo de Corriente de Falla Trifásica\\Sistema de Potencia II}
\author{TA Prof. Rafael Gil Bernal}
\date{27 de mayo de 2026}

\begin{document}
\maketitle

\section{Datos del Sistema}

\subsection{Topología}

El sistema tiene tres niveles de tensión y la siguiente configuración radial:

\[
\text{G} \xrightarrow[]{j0.6296} \boxed{A} \xrightarrow[T1]{j0.175} \boxed{B} \xrightarrow[L1+L2]{j0.42} \boxed{C} \xrightarrow[T2]{j0.2} \boxed{D}
\]

\[
\text{ Motor M }(j1.8) \parallel \text{ Carga CP }
\]

\subsection{Datos de los componentes (en pu, base 100 MVA)}

\begin{center}
\begin{tabular}{lccc}
\toprule
Componente & $S_{\text{nom}}$ & $X$ (base propia) & $X$ (base común) \\
\midrule
Generador G & 54 MVA & $X_d'' = 0.34$ pu & $j0.6296$ pu \\
T1 & 60 MVA & $X_{cc} = 10.5\%$ & $j0.175$ pu \\
L1 (20 km) & --- & $X_+ = 0.5$ $\Omega$/km & $j0.21$ pu \\
L2 (20 km) & --- & $X_+ = 0.5$ $\Omega$/km & $j0.21$ pu \\
T2 & 50 MVA & $X_{cc} = 10\%$ & $j0.2$ pu \\
Motor M & 10 MVA & $X_d'' = 0.18$ pu & $j1.8$ pu \\
\bottomrule
\end{tabular}
\end{center}

\subsection{Bases del sistema}

\[
S_{\text{base}} = 100\text{ MVA (trifásica)}
\]

\[
V_{\text{base},A} = V_{\text{base},D} = \frac{13.8}{\sqrt{3}}\text{ kV (fase-tierra)} \quad
V_{\text{base},B} = V_{\text{base},C} = \frac{69}{\sqrt{3}}\text{ kV (fase-tierra)}
\]

\[
I_{\text{base},A} = I_{\text{base},D} = \frac{S_{\text{base}}}{\sqrt{3}\,V_{LL}} = \frac{100\times10^3}{\sqrt{3}\times13.8} = 4183.7\text{ A}
\]

\[
I_{\text{base},B} = I_{\text{base},C} = \frac{100\times10^3}{\sqrt{3}\times69} = 836.7\text{ A}
\]

\[
Z_{\text{base},A} = Z_{\text{base},D} = \frac{V_{LL}^2}{S_{\text{base}}} = \frac{13.8^2}{100} = 1.9\;\Omega
\]

\[
Z_{\text{base},B} = Z_{\text{base},C} = \frac{69^2}{100} = 47.61\;\Omega
\]

\section{Casos de Pre-Falla}

\subsection{Caso base}
\begin{itemize}
    \item Tensión en barra D: $V_D = 0.98\angle0^\circ$ pu (fase-tierra)
    \item Motor: $S_M = 0.3$ pu, $\cos\phi = 0.95$ (atraso)
    \item Carga pasiva: $S_{CP} = 1.05$ pu, $\cos\phi = 0.85$ (atraso)
\end{itemize}

Corriente del motor (caso base):
\[
\dot{I}_M = \left(\frac{S_M}{V_D}\right)^* = \frac{0.3\angle-\cos^{-1}0.95}{0.98\angle0^\circ} = 0.102\angle-18.19^\circ\text{ pu}
\]

Corriente de la carga pasiva:
\[
\dot{I}_{CP} = \frac{1.05\angle-\cos^{-1}0.85}{0.98\angle0^\circ} = 0.3571\angle-31.79^\circ\text{ pu}
\]

Corriente total del generador:
\[
\dot{I}_G = \dot{I}_M + \dot{I}_{CP} = 0.102\angle-18.19^\circ + 0.3571\angle-31.79^\circ = 0.457\angle-28.78^\circ\text{ pu}
\]

\subsection{Caso máxima demanda}
\begin{itemize}
    \item Tensión en barra D: $V_D = 0.98\angle0^\circ$ pu
    \item Motor: $S_M = 0.6$ pu, $\cos\phi = 0.95$ (atraso)
    \item Carga pasiva: $S_{CP} = 1.05$ pu, $\cos\phi = 0.85$ (atraso)
\end{itemize}

Corriente del motor:
\[
\dot{I}_M = \frac{0.6\angle-18.19^\circ}{0.98\angle0^\circ} = 0.204\angle-18.19^\circ\text{ pu}
\]

Corriente total del generador:
\[
\dot{I}_G = 0.204\angle-18.19^\circ + 0.3571\angle-31.79^\circ = 0.552\angle-27.10^\circ\text{ pu}
\]

\section{Tensiones de Pre-Falla en Cada Barra}

Recorriendo el sistema desde la barra D hacia atrás (con $V_D = 0.98\angle0^\circ$ como referencia):

\subsection{Caso base}

\[
\dot{I}_G = 0.401 - j0.220\text{ pu}
\]

\textbf{Barra C:}
\[
\dot{V}_C = \dot{V}_D + jX_{T2}\,\dot{I}_G = 0.98 + j0.2(0.401 - j0.220)
\]
\[
\dot{V}_C = 0.98 + 0.044 + j0.0802 = 1.024 + j0.0802 = 1.027\angle4.48^\circ\text{ pu}
\]

\textbf{Barra B:}
\[
\dot{V}_B = \dot{V}_C + jX_{L12}\,\dot{I}_G,\quad X_{L12} = X_{L1} + X_{L2} = j0.42
\]
\[
\dot{V}_B = (1.024 + j0.0802) + 0.0924 + j0.1684 = 1.116 + j0.249 = 1.144\angle12.56^\circ\text{ pu}
\]

\textbf{Barra A:}
\[
\dot{V}_A = \dot{V}_B + jX_{T1}\,\dot{I}_G = (1.116 + j0.249) + 0.0385 + j0.0702
\]
\[
\dot{V}_A = 1.155 + j0.319 = 1.198\angle15.43^\circ\text{ pu}
\]

\textbf{Barra D:}
\[
\dot{V}_D = 0.98\angle0^\circ\text{ pu (dato)}
\]

\subsection{Caso máxima demanda}

\[
\dot{I}_G = 0.552\angle-27.10^\circ = 0.491 - j0.251\text{ pu}
\]

\textbf{Barra C:}
\[
\dot{V}_C = 0.98 + j0.2(0.491 - j0.251) = 0.98 + 0.050 + j0.098 = 1.030 + j0.098 = 1.035\angle5.44^\circ\text{ pu}
\]

\textbf{Barra B:}
\[
\dot{V}_B = (1.030 + j0.098) + j0.42(0.491 - j0.251) = 1.030 + j0.098 + 0.105 + j0.206
\]
\[
\dot{V}_B = 1.135 + j0.304 = 1.175\angle15.01^\circ\text{ pu}
\]

\textbf{Barra A:}
\[
\dot{V}_A = (1.135 + j0.304) + j0.175(0.491 - j0.251) = 1.135 + j0.304 + 0.044 + j0.086
\]
\[
\dot{V}_A = 1.179 + j0.390 = 1.242\angle18.31^\circ\text{ pu}
\]

\textbf{Barra D:}
\[
\dot{V}_D = 0.98\angle0^\circ\text{ pu (dato)}
\]

\section{Matriz de Admitancias (Y-Bus)}

Para la red de secuencia positiva con 4 barras:

Admitancias de cada elemento:
\[
y_G = \frac{1}{j0.6296} = -j1.588,\quad
y_{T1} = \frac{1}{j0.175} = -j5.714,\quad
y_{L12} = \frac{1}{j0.42} = -j2.381
\]
\[
y_{T2} = \frac{1}{j0.2} = -j5.0,\quad
y_M = \frac{1}{j1.8} = -j0.5556
\]

Matriz Y-Bus ($4\times4$):

\[
\mathbf{Y}_{\text{bus}} = 
\begin{bmatrix}
y_G + y_{T1} & -y_{T1} & 0 & 0\\[4pt]
-y_{T1} & y_{T1} + y_{L12} & -y_{L12} & 0\\[4pt]
0 & -y_{L12} & y_{L12} + y_{T2} & -y_{T2}\\[4pt]
0 & 0 & -y_{T2} & y_{T2} + y_M
\end{bmatrix}
\]

Sustituyendo:

\[
\boxed{
\mathbf{Y}_{\text{bus}} = 
\begin{bmatrix}
-j7.302 & j5.714 & 0 & 0\\[4pt]
j5.714 & -j8.095 & j2.381 & 0\\[4pt]
0 & j2.381 & -j7.381 & j5.0\\[4pt]
0 & 0 & j5.0 & -j5.5556
\end{bmatrix}
}
\]

\section{Matriz de Impedancias de Barra (Z-Bus)}

Para un sistema radial, la impedancia equivalente de Thevenin vista desde cada barra se obtiene combinando impedancias en serie y paralelo.

\subsection{Impedancia equivalente en Barra D ($Z_{44}$)}

Hacia el generador: $Z_{G} = j0.6296 + j0.175 + j0.42 + j0.2 = j1.4246$ pu

Hacia el motor: $Z_{M} = j1.8$ pu

\[
Z_{44} = \frac{j1.4246 \times j1.8}{j1.4246 + j1.8} = \frac{-2.564}{j3.2246} = j0.7953\text{ pu}
\]

\subsection{Impedancia equivalente en Barra C ($Z_{33}$)}

Hacia el generador: $j0.6296 + j0.175 + j0.21 = j1.0146$ pu (solo L1, sin L2)
Hacia el motor: $j0.2 + j1.8 = j2.0$ pu

\[
Z_{33} = \frac{j1.0146 \times j2.0}{j1.0146 + j2.0} = j0.6731\text{ pu}
\]

\subsection{Impedancia equivalente en Barra B ($Z_{22}$)}

Hacia el generador: $j0.6296 + j0.175 = j0.8046$ pu
Hacia el motor: $j0.21 + j0.21 + j0.2 + j1.8 = j2.42$ pu

\[
Z_{22} = \frac{j0.8046 \times j2.42}{j0.8046 + j2.42} = j0.6040\text{ pu}
\]

\subsection{Impedancia equivalente en Barra A ($Z_{11}$)}

Hacia el generador: $j0.6296$ pu
Hacia el motor: $j0.175 + j0.42 + j0.2 + j1.8 = j2.595$ pu

\[
Z_{11} = \frac{j0.6296 \times j2.595}{j0.6296 + j2.595} = j0.5067\text{ pu}
\]

\section{Corriente de Falla Trifásica Balanceada}

La corriente de falla trifásica en la barra $k$ es:

\[
I_{f,k} = \frac{V_k(0)}{Z_{kk}}
\]

\subsection{Caso base}

\begin{center}
\begin{tabular}{lccc}
\toprule
Barra & $V_k(0)$ (pu) & $Z_{kk}$ (pu) & $I_f$ (pu) \\
\midrule
A & $1.198\angle15.43^\circ$ & $j0.5067$ & $2.364\angle-74.57^\circ$ \\
B & $1.144\angle12.56^\circ$ & $j0.6040$ & $1.894\angle-77.44^\circ$ \\
C & $1.027\angle4.48^\circ$ & $j0.6731$ & $1.525\angle-85.52^\circ$ \\
D & $0.98\angle0^\circ$ & $j0.7953$ & $1.232\angle-90^\circ$ \\
\bottomrule
\end{tabular}
\end{center}

Valores reales (kA simétricos RMS):

\begin{center}
\begin{tabular}{lccc}
\toprule
Barra & $I_f$ (pu) & $I_{\text{base}}$ (A) & $I_f$ (kA) \\
\midrule
A & $2.364$ & 4183.7 & $9.89$ \\
B & $1.894$ & 836.7 & $1.58$ \\
C & $1.525$ & 836.7 & $1.28$ \\
D & $1.232$ & 4183.7 & $5.15$ \\
\bottomrule
\end{tabular}
\end{center}

\subsection{Caso máxima demanda}

\begin{center}
\begin{tabular}{lccc}
\toprule
Barra & $V_k(0)$ (pu) & $Z_{kk}$ (pu) & $I_f$ (pu) \\
\midrule
A & $1.242\angle18.31^\circ$ & $j0.5067$ & $2.451\angle-71.69^\circ$ \\
B & $1.175\angle15.01^\circ$ & $j0.6040$ & $1.945\angle-74.99^\circ$ \\
C & $1.035\angle5.44^\circ$ & $j0.6731$ & $1.538\angle-84.56^\circ$ \\
D & $0.98\angle0^\circ$ & $j0.7953$ & $1.232\angle-90^\circ$ \\
\bottomrule
\end{tabular}
\end{center}

Valores reales (kA simétricos RMS):

\begin{center}
\begin{tabular}{lccc}
\toprule
Barra & $I_f$ (pu) & $I_{\text{base}}$ (A) & $I_f$ (kA) \\
\midrule
A & $2.451$ & 4183.7 & $10.25$ \\
B & $1.945$ & 836.7 & $1.63$ \\
C & $1.538$ & 836.7 & $1.29$ \\
D & $1.232$ & 4183.7 & $5.15$ \\
\bottomrule
\end{tabular}
\end{center}

\section{Corriente Asimétrica Momentánea (Pico)}

La corriente asimétrica momentánea considera el peor caso de asimetría DC:

\[
I_{\text{mom}} = \sqrt{2} \times I_{\text{sym}} \times \text{FA}
\]

donde FA es el factor de asimetría:
\[
\text{FA} = 1 + e^{-\pi/(X/R)}
\]

Sin datos de resistencia, se asume un factor conservador $\text{FA} = 1.8$ (corresponde a $X/R \approx 15$).

\subsection{Caso base}

\begin{center}
\begin{tabular}{lccc}
\toprule
Barra & $I_{\text{sym}}$ (kA RMS) & $I_{\text{mom}}$ (kA pico) \\
\midrule
A & $9.89$ & $25.18$ \\
B & $1.58$ & $4.02$ \\
C & $1.28$ & $3.26$ \\
D & $5.15$ & $13.11$ \\
\bottomrule
\end{tabular}
\end{center}

\subsection{Caso máxima demanda}

\begin{center}
\begin{tabular}{lccc}
\toprule
Barra & $I_{\text{sym}}$ (kA RMS) & $I_{\text{mom}}$ (kA pico) \\
\midrule
A & $10.25$ & $26.09$ \\
B & $1.63$ & $4.15$ \\
C & $1.29$ & $3.28$ \\
D & $5.15$ & $13.11$ \\
\bottomrule
\end{tabular}
\end{center}

\section{Corriente Asimétrica de Interrupción (6 ciclos)}

A 60 Hz, 6 ciclos corresponden a $t = 0.1$ s. Para este tiempo de interrupción:

\[
I_{\text{DC}}(t) = \sqrt{2}\,I_{\text{sym}}\,e^{-\omega t/(X/R)}
\]

\[
I_{\text{asym,int}}(t) = \sqrt{I_{\text{sym}}^2 + I_{\text{DC}}(t)^2}
\]

Para $X/R = 15$ (valor típico asumido):
\[
\frac{I_{\text{DC}}}{I_{\text{sym}}\sqrt{2}} = e^{-2\pi\times60\times0.1/15} = e^{-2.513} = 0.081
\]

\[
I_{\text{asym,int}} = I_{\text{sym}}\sqrt{1 + 2(0.081)^2} = 1.007\,I_{\text{sym}}
\]

\subsection{Caso base}

\begin{center}
\begin{tabular}{lccc}
\toprule
Barra & $I_{\text{sym}}$ (kA) & $I_{\text{asym,int}}$ (kA) \\
\midrule
A & $9.89$ & $9.96$ \\
B & $1.58$ & $1.59$ \\
C & $1.28$ & $1.29$ \\
D & $5.15$ & $5.19$ \\
\bottomrule
\end{tabular}
\end{center}

\subsection{Caso máxima demanda}

\begin{center}
\begin{tabular}{lccc}
\toprule
Barra & $I_{\text{sym}}$ (kA) & $I_{\text{asym,int}}$ (kA) \\
\midrule
A & $10.25$ & $10.32$ \\
B & $1.63$ & $1.64$ \\
C & $1.29$ & $1.30$ \\
D & $5.15$ & $5.19$ \\
\bottomrule
\end{tabular}
\end{center}

\section{Resumen de Resultados}

\subsection{Caso base}

\begin{center}
\begin{tabular}{lcccc}
\toprule
Barra & Tensión pre-falla (pu) & $I_{\text{sym}}$ (kA) & $I_{\text{mom}}$ (kA pico) & $I_{\text{int,6cic}}$ (kA) \\
\midrule
A & $1.198\angle15.43^\circ$ & 9.89 & 25.18 & 9.96 \\
B & $1.144\angle12.56^\circ$ & 1.58 & 4.02 & 1.59 \\
C & $1.027\angle4.48^\circ$ & 1.28 & 3.26 & 1.29 \\
D & $0.98\angle0^\circ$ & 5.15 & 13.11 & 5.19 \\
\bottomrule
\end{tabular}
\end{center}

\subsection{Caso máxima demanda}

\begin{center}
\begin{tabular}{lcccc}
\toprule
Barra & Tensión pre-falla (pu) & $I_{\text{sym}}$ (kA) & $I_{\text{mom}}$ (kA pico) & $I_{\text{int,6cic}}$ (kA) \\
\midrule
A & $1.242\angle18.31^\circ$ & 10.25 & 26.09 & 10.32 \\
B & $1.175\angle15.01^\circ$ & 1.63 & 4.15 & 1.64 \\
C & $1.035\angle5.44^\circ$ & 1.29 & 3.28 & 1.30 \\
D & $0.98\angle0^\circ$ & 5.15 & 13.11 & 5.19 \\
\bottomrule
\end{tabular}
\end{center}

\section*{Notas}
\begin{enumerate}
    \item Los valores de resistencia no fueron proporcionados; se asumió $X/R = 15$ para el cálculo de la componente asimétrica.
    \item La corriente momentánea usa factor de asimetría $\text{FA} = 1.8$.
    \item La corriente de interrupción a 6 ciclos considera la atenuación de la componente DC con constante de tiempo $\tau = X/(\omega R)$.
    \item Las tensiones de pre-falla se calcularon recorriendo el sistema desde la barra D (referencia $0.98\angle0^\circ$ pu) hacia el generador.
    \item En la barra D, la corriente de falla es la misma en ambos casos porque la impedancia de Thevenin no cambia y la tensión pre-falla es la misma (dato del problema).
\end{enumerate}

\end{document}
